\documentclass[../../../include/open-logic-section]{subfiles}

\begin{document}

\olfileid{sth}{z}{pairs}
\olsection{జంటలు}

ఇప్పుడు పరిశీలించాల్సిన
స్వీకృతం ఇది:

\begin{axiom}[జంటలు]
ఏ సమితులు $a, b$కైనా,
$\{a, b\}$ సమితి ఉంటుంది.
\[
	\forall a \forall b \exists P \forall x (x \in P \liff (x = a \lor x = b))
\]
\end{axiom}

దశలవారీ సమితి భావనతో
ఈ స్వీకృతాన్ని ఇలా
సమర్థించవచ్చు. $a$ సమితి
$S$ దశలో, $b$ సమితి
$T$ దశలో అందుబాటులో
ఉన్నాయనుకుందాం. $S$,
$T$ దశల్లో తరువాత
వచ్చేదాన్ని $M$ అనుకుందాం.
అప్పుడు $a$, $b$ రెండూ
$M$ దశలో అందుబాటులో
ఉన్నాయి; కనుక $\{a,b\}$
అనే సమితిని $M$ తరువాతి
ఏ దశలోనైనా ఏర్పరచవచ్చు.

కానీ ఆగండి!{} $M$ తరువాత
ఏవైనా దశలు \emph{ఉన్నాయని}
ఎందుకు అనుకోవాలి? అవి
లేకపోతే మన సమర్థన
విఫలమవుతుంది. కాబట్టి
జంటల స్వీకృతాన్ని
సమర్థించాలంటే,
\olref[sth][z][story]{sec}లోని
కథకు మరో సూత్రం
జోడించాలి. అది:
\begin{enumerate}
	\item[] \stagessucc. చివరి దశ లేదు.
\end{enumerate}
ఈ సూత్రానికి సమర్థన ఉందా?
చివరి దశ లేదని
షోన్‌ఫీల్డ్ కథ \emph{స్పష్టంగా}
చెప్పలేదు. అయినా ఇది
మన కథకు (కచ్చితంగా
చెప్పాలంటే) అదనపు సూత్రమే
అయినప్పటికీ, సమితులు
దశలవారీగా ఏర్పడతాయనే
మూల భావనకు సరిపోతుంది.
ఇకముందు దీనిని అంగీకరిస్తాం.
దానితో పాటు జంటల
స్వీకృతాన్నీ అంగీకరిస్తాం.

ఈ కొత్త స్వీకృతంతో
మరికొన్ని సమితులు
ఉంటాయని నిరూపించవచ్చు.
ఉదాహరణకు:

\begin{prop}\ollabel{prop:pairsconsequences}
ఏ సమితులు $a$, $b$కైనా,
ఈ కింది సమితులు ఉంటాయి:
	\begin{enumerate}
		\item\ollabel{singleton} $\{a\}$
		\item\ollabel{binunion} $a \cup b$
		\item\ollabel{tuples} $\tuple{a, b}$ క్రమయుగ్మం
	\end{enumerate}
\end{prop}

\begin{proof}
\olref{singleton}. జంటల స్వీకృతం
ప్రకారం $\{a, a\}$ ఉంటుంది;
సమితుల సమానత్వ సూత్రం
ప్రకారం అది $\{a\}$.

\olref{binunion}. జంటల
స్వీకృతం ప్రకారం
$\{a, b\}$ ఉంటుంది.
ఇప్పుడు సమ్మేళన స్వీకృతం
ప్రకారం $a \cup b = \bigcup
\{a, b\}$ ఉంటుంది.

\olref{tuples}. \olref{singleton}
ప్రకారం $\{a\}$ ఉంటుంది.
జంటల స్వీకృతం ప్రకారం
$\{a,
b\}$ ఉంటుంది. మళ్లీ
జంటల స్వీకృతం ప్రకారం
$\{\{a\}, \{a, b\}\} = \tuple{a, b}$
ఉంటుంది.
\end{proof}

\begin{prob}
ఏ సమితులు $a, b, c$కైనా,
$\{a, b, c\}$ సమితి
ఉంటుందని చూపండి.
\end{prob}

\begin{prob}
ఏ సమితులు $a_1, \ldots, a_n$కైనా,
$\{a_1, \ldots,
a_n\}$ సమితి ఉంటుందని
చూపండి.
\end{prob}

%\begin{proof} జంటల స్వీకృతాన్ని రెండుసార్లు ప్రయోగిస్తే, $\{\{a_1, a_2\}, \{a_1,
%   a_3\}\}$ ఉంటుంది. సమ్మేళనం, సమితుల సమానత్వ సూత్రం ప్రకారం, $\bigcup \{\{a_1,
%   a_2\}, \{a_1, a_3\}\} = \{a_1, a_2, a_3\}$ ఉంటుంది. అవసరమైనన్ని
%   సార్లు ఈ పద్ధతిని పునరావృతం చేయండి. \end{proof}

\end{document}
